% Только для преподавателя. Студент защищает один выбранный вариант работы.
\providecommand{\defenseq}[2]{\par\medskip\noindent\textbf{#1}\par\textit{Признаки хорошего ответа:} #2\par}

\section*{Часть I. Краткий маршрут защиты}
Начните с отчёта и работающего результата. Выберите два общих вопроса и
вопросы \emph{только по выполненному студентом варианту}. Попросите показать
свои команды, исходники, ожидаемый и наблюдаемый результат. Устная защита
следует после короткого письменного теста; тест не является допуском.

\subsection*{Общие вопросы}
\defenseq{Где находилась программа до запуска и откуда CPU получает её инструкции во время работы?}{Различает файл на накопителе и код в RAM; называет CPU исполнителем инструкций, а не постоянным хранилищем. Просит показать исходник и результат запуска.}
\defenseq{Чем адрес отличается от значения по этому адресу?}{На примере собственного массива или переменной показывает место хранения и содержимое; различает регистр, RAM и накопитель.}

\subsection*{Вариант 1: факториал на C}
\defenseq{Почему после изменения \texttt{factorial.c} нужно повторить команду GCC?}{Различает исходник и исполняемый файл; показывает сборку и запуск \texttt{./factorial}, объясняет прежний результат без пересборки.}
\defenseq{Как в GDB вы нашли и исправили ошибку \texttt{i < n}?}{Показывает собственную остановку и значения \texttt{i/result}; при \texttt{n=5} пропущен множитель 5, получается 24 вместо 120; исправляет и пересобирает.}

\subsection*{Вариант 2: NASM, BIOS и QEMU}
\defenseq{Покажите путь оценки от клавиши до сохранения.}{Различает скан-код, числовой байт в RAM, символ и атрибут VGA, секторный буфер и сектор отдельного HDD. Показывает свой изменённый \texttt{journal.asm} и экран/байты образа.}
\defenseq{Почему при изменении допустимой оценки и сектора сохранения пришлось исправить несколько мест?}{Показывает ввод значения, проверку при загрузке, обновлённый DAP и проверку байтов по правильному смещению; пересчитывает связанные значения, не полагаясь только на сообщение \texttt{SAVED}.}

\section*{Часть II. Банк дополнительных вопросов}
Выбирайте вопросы по тому, что студент сделал и зафиксировал в отчёте.
Признаки служат ориентирами для обсуждения, а не текстом для заучивания.

\subsection*{Общая архитектура и работа с памятью}
\defenseq{Что делают CPU, RAM и SSD/HDD в вашей работе?}{CPU исполняет инструкции, RAM хранит рабочие данные и код во время исполнения, накопитель сохраняет без питания; показывает пример из своего проекта.}
\defenseq{Что означает адресация по байтам?}{Каждый байт имеет адрес; соседние байты массива следуют подряд; не смешивает адрес с записанным в байт значением.}
\defenseq{Чем регистр отличается от ячейки RAM?}{Регистр находится внутри CPU и используется операциями напрямую; RAM адресуется отдельно; регистр не заменяет всё адресное пространство.}
\defenseq{Почему \texttt{sizeof} и адрес конкретной переменной надо измерять на выбранной платформе?}{Размер зависит от типа и платформы, адрес от запуска и окружения; не объявляет адрес или разрядность \texttt{long} универсальными.}
\defenseq{Что хранится в переменной-указателе и что даёт её разыменование?}{Указатель содержит адрес; разыменование обращается к объекту по нему; показывает адрес и значение на своём примере.}
\defenseq{Почему элементы массива удобно обходить последовательно?}{Их размещение позволяет вычислить смещение по индексу; соседние элементы могут использовать одну кэш-линию; не обещает конкретную скорость в QEMU.}
\defenseq{Зачем нужен стек при вызове функции?}{Объясняет адрес возврата, кадры вызовов и локальные данные; показывает цепочку вызовов или сохранение регистров.}
\defenseq{Чем стек отличается от кучи на уровне назначения?}{Стек связан с вызовами, куча выделяется динамически; не приписывает автоматическое освобождение памяти кучи возврату из функции.}
\defenseq{Чем код в RAM отличается от данных в RAM?}{Оба представлены байтами, но CPU выбирает и декодирует инструкции по адресу исполнения, а операнды читает или изменяет отдельно.}
\defenseq{Зачем компьютеру кэш, если уже есть RAM?}{Кэш хранит копии часто используемых линий ближе к CPU; промах требует обращения к более медленной памяти; результат не равен постоянному хранению.}
\defenseq{Когда устройство может потребовать внимания CPU?}{Сравнивает опрос готовности и аппаратное прерывание; устройство уведомляет, а программа затем обрабатывает событие.}
\defenseq{Чем BIOS-прошивка отличается от операционной системы?}{BIOS помогает начать загрузку и даёт аппаратные сервисы; ОС предоставляет процессы, файлы и управление ресурсами; QEMU-журнал работает без гостевой ОС.}

\subsection*{Вариант 1: C, GCC и GDB}
\defenseq{Что именно означают \texttt{factorial.c}, \texttt{factorial} и запущенная программа?}{Исходный текст, готовый исполняемый файл и процесс --- разные объекты; показывает каждый в своей рабочей папке.}
\defenseq{За что отвечают \texttt{-Wall}, \texttt{-g}, \texttt{-O0} и \texttt{-o}?}{Предупреждения, отладочные сведения, отключение оптимизации для понятных шагов и имя выходного файла; не путает \texttt{-g} с запуском GDB.}
\defenseq{Почему \texttt{0! = 1} в вашей функции?}{Результат инициализирован единицей, цикл для нуля не выполняет умножений; показывает запуск с нулём.}
\defenseq{Зачем в работе ограничение входа 0--20?}{Отрицательный факториал здесь не предусмотрен; \texttt{20!} помещается в \texttt{unsigned long long}, \texttt{21!} выходит за его диапазон на учебной платформе.}
\defenseq{Как обнаруживается неверный ввод?}{Проверяет возвращённое значение \texttt{fscanf} и диапазон \texttt{n}; выводит сообщение об ошибке и ненулевой код выхода.}
\defenseq{Куда идут \texttt{stdin}, \texttt{stdout} и \texttt{stderr} при перенаправлении?}{Вход берётся из файла, обычный результат и диагностика могут идти в разные файлы; наличие текста в \texttt{stderr} само по себе не означает неуспех.}
\defenseq{Что показывает \texttt{echo \$?} и когда его надо вызывать?}{Код завершения последней запущенной команды; проверяет сразу после программы, не после промежуточного \texttt{ls} или \texttt{cat}.}
\defenseq{Чем \texttt{break} отличается от \texttt{watch}?}{Первое останавливает у строки или функции, второе --- при изменении значения; студент показывает наблюдение за \texttt{result}.}
\defenseq{Чем различаются \texttt{next}, \texttt{step} и \texttt{finish}?}{Шаг без входа в функцию, вход в вызов и продолжение до возврата; один \texttt{next} не обязан равняться одной итерации цикла.}
\defenseq{Что показывают \texttt{backtrace} и \texttt{frame}?}{Цепочку вызовов и выбранный кадр; локальные переменные доступны лишь в соответствующем контексте.}
\defenseq{Как прочитать байты переменной \texttt{result} в GDB?}{Показывает \texttt{sizeof(result)}, \texttt{\&result} и \texttt{x/8xb}; при 120 первым идёт байт \texttt{78h} на учебной x86-64.}
\defenseq{Почему ошибка \texttt{i < n} даёт 24 для \texttt{n=5}?}{Цикл умножает только на 1, 2, 3, 4; объясняет, как нашёл границу, вернул \texttt{i <= n} и повторно собрал программу.}

\subsection*{Вариант 2: NASM, BIOS и QEMU}
\defenseq{Как BIOS находит и запускает ваш первый сектор?}{Сектор занимает 512 байтов, подпись \texttt{55 AA} находится в конце, BIOS помещает его в RAM по адресу \texttt{0x7C00} и передаёт управление.}
\defenseq{Зачем после NASM создавать образ дискеты?}{\texttt{-f bin} даёт плоские байты, образ включает загрузчик и журнал; показывает три операции \texttt{dd} и не перезаписывает диск оценок.}
\defenseq{Почему в QEMU два образа дисков?}{На дискете код для загрузки, на отдельном HDD сектор оценок; имена файлов на хосте не означают файловую систему внутри гостя.}
\defenseq{Как вычислить позицию символа на текстовом экране?}{Одна ячейка содержит код символа и атрибут; 80 столбцов дают 160 байтов на строку; смещение равно \texttt{160*строка+2*столбец}.}
\defenseq{Как формируется байт цвета?}{Младшие четыре бита --- цвет символа, старшие --- фон/мигание; студент показывает собственный расчёт и меняет нужный операнд \texttt{vga\_text}, а не очистку экрана.}
\defenseq{Почему для Q/E нельзя подставлять ASCII-коды букв вместо скан-кодов?}{Обработчик сопоставляет события клавиатуры со скан-кодами; переводит данные числа в hex и показывает две изменённые ветви кода и подсказку на экране.}
\defenseq{Чем вызов \texttt{INT 16h} отличается от IRQ1?}{Первый инициирует инструкция программы, второй сообщает аппаратное событие от клавиатуры; после установки своего обработчика стандартная очередь BIOS не пополняется.}
\defenseq{Как представлены 16 студентов с четырьмя оценками?}{64 последовательных байта; адрес поля зависит от \texttt{4*i+j}; \texttt{FFh} означает отсутствие оценки, ноль --- настоящая выставленная оценка.}
\defenseq{Как рассчитать среднее с пропусками и нулём?}{Суммирует и считает только выставленные оценки, ноль включает в число оценок, \texttt{FFh} пропускает; при пустом наборе не делит на ноль.}
\defenseq{Что изменяется при расширении шкалы до оценки 12?}{Ввод и проверка сохранённых оценок должны согласоваться; показывает среднее \texttt{[12,0]} как 6.00 и \texttt{[12,1,1]} как 4.67.}
\defenseq{Как перенести данные на другой сектор диска?}{В DAP меняет LBA; для смещения 4096 байтов и сектора 512 байтов получает LBA 8; проверяет \texttt{GR16} именно там.}
\defenseq{Почему повреждённый сектор не должен затереть рабочую таблицу?}{Сначала сверяются сигнатура, версия, контрольная сумма и допустимые оценки; только после проверки 64 байта копируются в RAM; показывает \texttt{BAD\_CHECKSUM} на отдельной копии.}
